\documentclass{article}
\usepackage{../math_template}

\author{Jonathan Kasongo}
\title{Codeforces Round 1075 --- Problems C1 and C2}

\begin{document}
\maketitle
\textit{The problem statement is worded differently here.}
\begin{problem}{Codeforces Round 1075 --- Problems C1 and C2}
Given a natural number $n$. Find a permutation $p$ of $(1,2,\ldots, n)$,
such that for each $i \in S$ there exists a $i \leq j \leq n$
such that $p_i = p_j \oplus i$.\vspace{0.2cm}

C1 --- Solve the problem for $S = \{2, 3, \ldots, n-1\}$.\vspace{0.2cm}

C2 --- Solve the problem for $S = \{ 1, 2, \ldots, n-1\}$.\\

\textbf{Input}

Each test contains multiple test cases. The first line contains the number
of test cases $1 \leq t \leq 10^4$.\vspace{0.2cm}

The description of the test cases follows.
The only line of each test case contains one positive integer
$3 \leq n \leq 2 \cdot 10^5$, the length of the permutation. \vspace{0.2cm}

It is guaranteed that the sum of $n$ over all test cases does not exceed
$2 \cdot 10^5$.\\

\textbf{Output}

If there is a suitable permutation output the $n$ integers
$p_1 \ p_2 \ ...\ p_n$ with spaces between them. Otherwise output $-1$.
\end{problem}

\begin{solution}{Solution}
\textit{If we leave out any details in the solution then it's because,
either it shouldn't be too difficult to fill in the gaps.} \vspace{0.2cm}

C1 --- We look for an easy way to satisfy the desired condition. After some
exploration one may notice that
$$
a = \begin{cases}
(a-1) \oplus 1 \quad \text{if } a \text{ is odd,}\\
(a+1) \oplus 1 \quad \text{if } a \text{ is even.}
\end{cases}
$$
This is the heart of the problem.
The claim is true because if $a$ is odd then $(a-1)_2$ is just $(a)_2$
without the $2^0$ bit. Now if you XOR with 1, you are adding the $2^0$
bit back in, giving $a = (a-1) \oplus 1$. Similarly, if $a$ is even then
$(a+1)_2$ is just $(a)_2$ with a $2^0$ bit added.$^{*}$  Thus if one XOR's
with 1 then you will remove this $2^0$
bit again giving $a = (a+1) \oplus 1$. \vspace{0.2cm}

So if we place $1$ at the end of our permutation and let
$$
p_i = \begin{cases}
i+1 \quad \text{if } i \text{ is even,}\\
i-1 \quad \text{if } i \text{ is odd.}
\end{cases}
$$
then just let $p_1$ take on the one unused value,
then we can take $j = n$ to get $p_n \oplus i = p_i$, and this can be
checked to work for all suitable $2 \leq i \leq n-1$.
\vspace{0.2cm}

\scriptsize{* Note that $(a)_2$ didn't have
this bit before because only odd numbers have this bit in their binary
representation.}\\

\normalsize
C2 --- One can check that if $n$ is odd then our construction from C1 works
fine. The only issue is with even $n$. Again after some experimentation one
may find that the construction is never possible if $n$ is a power of 2.
\vspace{0.2cm}

This is true because if $a \oplus b = 2^k$ for $1 \leq b \leq a \leq 2^k$
then $a$ must have the $2^k$ bit in it's binary representation. Now we
notice that whatever bits $b$ has (other than the $2^k$ bit), $a$ must also
have these bits to satisfy the XOR equation. Since $b \geq 1$ we have that
$a \geq 2^k + 1$ a contradiction.\vspace{0.2cm}

Now if $n$ is even but isn't a power of 2, then we look for a construction
that works in this occasion. Let's start out from our construction in
C1. If you write out some cases, you'll find that
the only issue occurs when $i = 1$ and we want to find $j$ such that
$p_1 = n = p_j \oplus 1$.
It is natural then to guess that perhaps we can use
our construction from C1, except just swap $p_1$ and some other number.
Let $l \geq 0$ be an integer such that $2^l$ is the
smallest bit that occurs in the binary representation of $n$.
If you swap the elements at $p_1$ and $p_{2^l}$, then the resultant
permutation satisfies the desired criteria. \textit{We leave the proof
of this fact as an exercise to the reader, as a hint look at the number
$n - 2^l$.} $\blacksquare$
\end{solution}

\end{document}
